%% Согласно ГОСТ Р 7.0.11-2011:
%% 5.3.3 В заключении диссертации излагают итоги выполненного исследования, рекомендации, перспективы дальнейшей разработки темы.
%% 9.2.3 В заключении автореферата диссертации излагают итоги данного исследования, рекомендации и перспективы дальнейшей разработки темы.


\textit{В заключении подводятся итоги работы.
Формулируются основные выводы (7-9) по результатам исследований.
Приводятся сведения об апробации, полноте опубликования в научной печати основного содержания диссертации, ее результатов, выводов.
Приводятся сведения о защищенности технических решений авторскими свидетельствами (патентами).
Указываются предприятия, где внедрены результаты диссертационной работы и где еще они могут быть использованы.
Этот раздел занимает до восьми страниц текста.
Можно построить заключение к диссертации по схеме выполнения общей характеристики работы, приводимой в автореферате, что позволит усилить единство диссертации и автореферата.}

В диссертационной работе решена актуальная научно-техническая задача повышения коэффициента напора осевых вентиляторов при сохранении их энергетической эффективности за счёт применения меридионального ускорения потока и разработки методов их проектирования.

\begin{enumerate}
  \item Разработана и реализована математическая модель квазитрёхмерного течения в лопаточных венцах вентиляторов с меридиональным ускорением, основанная на модификации метода дискретных вихрей и конформном отображении осесимметричных поверхностей тока переменной толщины.
  Результатом численного моделирования является поле скоростей и давления в межлопаточном канале вращающихся и неподвижных пространственных решёток профилей.

  Основное отличие от существующих решений заключается в том, что модель основана на особенностях потенциального течения, что определяет её преимущества:
 \begin{enumerate}
  \item не требуется построения расчётной сетки и задания начального приближения;

  \item моделирование сводится к решению СЛАУ с относительно небольшим количеством неизвестных (примерно 150);

  \item результаты моделирования при различных режимах течения являются линейнозависимыми, течение на любом режиме работы решётки профилей моделируется как линейная комбинация результатов двух известных режимов.
\end{enumerate}

  \item На базе модели течения разработан метод профилирования лопаток вентилятора с меридиональным ускорением потока, осевых, диагональных и центробежных турбомашин.
  Использование теоретических характеристик решёток профилей позволило распространить подходы к профилированию осевых турбомашин на машины с произвольными меридиональными формами.
  Проведённые проверочные RANS-расчёты продемонстрировали удовлетворительное совпадение распределения давления на контурах лопаток, а так же расчётных коэффициентов теоретического напора в пределах 3,5\% до уточнения влияния вязкости, и менее 1\% после уточнения.

  \item Проведённое экспериментальное исследование характеристик модели вентилятора с меридиональным ускорением, спрофилированного разработанным методом, позволило подтвердить расчётные интегральные характеристики, такие как коэффициент теоретического напора \(\bar{H}_\text{т}\), и распределения \(\bar{H}_\text{т}\) и меридиональных компонент скорости \(\bar{c}_2m\) по высоте проточной части.
  Разница между расчётным и экспериментальным значениями \(\bar{H}_\text{т}\) составила менее 1\% на номинальном режиме и менее 3\% на остальной исследованной части рабочей характеристики.

  \item С использованием разработанных методов математического моделирования течения, было исследование влияние формы проточной части на характеристики вентиляторов с меридиональным ускорением, подтверждённые RANS-расчётами и экспериментом:
  \begin{enumerate}
  \item Показана зависимость коэффициента теоретического напора \(\bar{H}_\text{т}\) от сужения проточной части и  предложен способ оценки его влияния.
  В сравнении с экспериментом ошибка оценки \(\bar{H}_\text{т}\) составила менее 5\% при расширении проточной части и менее 3\% при сужении;

  \item Показано влияние меридионального ускорения на возможность повышения коэффициента теоретического напора \(\bar{H}_\text{т}\) для элементарного рабочего колеса, элементарной и полной ступени.
  Показано, что при сохранении КПД ступени, \(\bar{H}_\text{т}\) может быть значительно увеличен.
  В зависимости от относительного диаметра втулки и отношения динамического напора ступени к полному, \(\bar{H}_\text{т}\) может быть увеличен более чем в два раза, в сравнении с цилиндрической ступенью.
  Это позволяет при фиксированных значениях напора, расхода и мощности снизить частоту вращения ротора на 30\%, что приводит к снижению уровня суммарной звуковой мощности на 4,65~дБ (в 2,92 раз).
  \end{enumerate}

\end{enumerate}
